\section{Marco teórico y antecedentes}

\subsection{Bases Teóricas}
Esta sección fundamenta los conceptos pilares sobre los cuales se construye el mecanismo
de detección automatizada.
\begin{itemize}
    \item \textbf{Deriva y Erosión Arquitectónica:} Se define como la
        divergencia progresiva entre la arquitectura prescrita (planeada) y la
        arquitectura realmente implementada en el código fuente. Este fenómeno implica
        una pérdida de integridad estructural debido a la violación de principios
        de diseño originales \parencite{Pigazzini2021}. La investigación reciente
        destaca que esta divergencia es un problema persistente cuando no existen
        acciones de contención, impactando negativamente en la mantenibilidad y
        la capacidad de evolución del sistema \parencite{Anthony2024}. Se estima
        que el costo del mantenimiento de software, a menudo inflado por la erosión,
        consume aproximadamente el 65\% del presupuesto de TI de una organización
        \parencite{Pigazzini2021}.

    \item \textbf{Deuda Técnica Arquitectónica (ATD):} Es una metáfora que
        describe las decisiones de diseño subóptimas que ofrecen beneficios a
        corto plazo (como entregas rápidas) pero deterioran la calidad a largo
        plazo, generando un "interés" en forma de mayor esfuerzo de mantenimiento
        \parencite{Hamed2023}. La ATD se diferencia de la deuda a nivel de código
        porque involucra decisiones estructurales globales como capas,
        subsistemas e interfaces \parencite{Pigazzini2021}. Un síntoma clave de
        su acumulación es la presencia de Smells Arquitectónicos (AS) \parencite{SasAvgeriou}.

    \item \textbf{Smells Arquitectónicos (AS):} Son decisiones de diseño que
        impactan negativamente en la calidad interna del sistema \parencite{Pigazzini2021}.
        A diferencia de los code smells, que ocurren a nivel de métodos o clases,
        los AS involucran múltiples componentes, paquetes o capas \parencite{SasAvgeriou}.
        Ejemplos críticos incluyen la \textbf{Dependencia Cíclica (CD)}, donde los
        componentes no pueden ser liberados o probados de forma aislada, y el
        \textbf{Componente Dios (GC)}, que ocurre cuando una entidad es
        excesivamente grande y difícil de entender \parencite{SasAvgeriou, Pigazzini2021}.

    \item \textbf{Funciones de Aptitud (Fitness Functions):} Una función de
        aptitud es un mecanismo diseñado para resumir qué tan cerca está una
        solución de diseño específica de alcanzar sus objetivos y preservar sus
        características clave. Se definen formalmente como "cualquier mecanismo
        que proporciona una evaluación de integridad objetiva de alguna
        característica arquitectónica". Su propósito fundamental es actuar como
        barandillas (guardrails) automatizadas que comunican, validan y preservan
        atributos de calidad (requisitos no funcionales) de manera continua,
        asegurando que la evolución del sistema se mantenga dentro de un rango y
        una dirección deseados por los arquitectos. Esta disciplina permite a los
        equipos ser reactivos ante la degradación de características críticas en
        tiempo real, evitando que la realidad del sistema en producción se
        desvíe de las intenciones de diseño originales \parencite{FitnessF}.

    \item \textbf{Análisis Estático y Grafos de Dependencia:} Es la técnica de
        extraer hechos y abstracciones del código fuente sin ejecutarlo
        \parencite{Pigazzini2021}. Los grafos de dependencia representan los
        componentes del sistema como nodos y sus relaciones como aristas, permitiendo
        aplicar algoritmos de búsqueda (como \textit{Depth First Search}) para
        detectar estructuras problemáticas como ciclos \parencite{Pigazzini2021, Klein2022}.

    \item \textbf{Algoritmos de Recorrido sobre el Grafo de Dependencias:} La
        detección de smells arquitectónicos sobre el grafo $G=(V,E)$ requiere la
        aplicación de algoritmos de grafos dirigidos. La \textbf{Dependencia
        Cíclica (CD)} se detecta mediante una búsqueda en profundidad (\textit{DFS})
        que identifica Componentes Fuertemente Conexos (SCC), clasificando un
        subconjunto de nodos como cíclico cuando $|\text{SCC}| > 1$. La \textbf{Dependencia
        Inestable (UD)} se calcula midiendo el acoplamiento aferente ($C_{a}$ = aristas
        entrantes) y eferente ($C_{e}$ = aristas salientes) de cada nodo, derivando
        la métrica $I = C_{e}\,/\,(C_{a} + C_{e})$; una relación $A \to B$
        constituye un UD cuando el componente estable $A$ depende del inestable
        $B$ ($I_{A} < I_{B}$). El \textbf{Componente Dios (GC)} se detecta
        cuando el número de dependencias de salida de un paquete supera un
        umbral configurable. \textcite{Sas2019} validaron empíricamente estas
        métricas analizando la evolución de smells de inestabilidad en proyectos
        Java de código abierto mediante el recorrido sistemático del grafo de
        dependencias a nivel de paquete \parencite{Sas2019}.
\end{itemize}

\subsection{Antecedentes}
La investigación actual presenta diversas estrategias para gestionar la integridad
técnica del software, evolucionando desde la medición manual hacia enfoques
predictivos y automatizados.
\begin{itemize}
    \item \textbf{ArchUnit:} Es una biblioteca gratuita y extensible diseñada
        para verificar la arquitectura de código Java mediante cualquier entorno
        de pruebas unitarias estándar. Funciona analizando el bytecode e importando
        las clases en una estructura de código Java, lo que permite comprobar
        dependencias entre paquetes, capas, ciclos y más. Es la herramienta que permite
        materializar el concepto de "arquitectura auto-testeable" al permitir escribir
        reglas arquitectónicas como si fueran pruebas unitarias comunes \parencite{ArchUnit}.

    \item \textbf{ArchSync:} Este es una herramienta de asistencia que busca
        mantener en sincronía el código con los escenarios arquitectónicos (Use Case
        Maps). ¿Qué hace?, Utiliza heurísticas para procesar trazas de ejecución
        y correlacionarlas con el comportamiento arquitectónico esperado,
        sugiriendo "scripts de actualización" para corregir inconsistencias. Aunque
        ArchSync se apoya más en el comportamiento dinámico (trazas), coincide
        en reducir el esfuerzo de sincronización manual entre arquitectos y
        desarrolladores. \parencite{ArchSyncJAndres}

    \item \textbf{Gestión en PyMES:} \parencite{Hamed2023} proponen un enfoque de
        gestión de deuda técnica diseñado específicamente para pequeñas y medianas
        empresas (SMEs). Destacan que las limitaciones de recursos, presupuestos
        ajustados y la falta de documentación actualizada en estos entornos hacen
        vital una gobernanza basada en herramientas de análisis estático (como
        SonarQube) que operen sobre el artefacto vivo del software \parencite{Hamed2023}.

    \item \textbf{Índices de Deuda Basados en Machine Learning:} \parencite{SasAvgeriou}
        desarrollaron el \textbf{ATDI (Architectural Technical Debt Index)}, el cual
        utiliza modelos de aprendizaje automático (gradient boosting) para clasificar
        la severidad de los smells arquitectónicos. Este enfoque combina el
        análisis estático de las líneas de código afectadas con la severidad
        estimada, logrando un 71\% de acuerdo con la percepción de esfuerzo de
        refactorización de desarrolladores expertos \parencite{SasAvgeriou}.

    \item \textbf{Conformidad Automatizada en el Flujo de CI/CD:} \parencite{Klein2022}
        subrayan que el uso de verificadores de conformidad automatizados integrados
        en los flujos de Integración Continua (CI) permite exponer las desviaciones
        en el momento del commit. Esto facilita la remediación inmediata antes
        de que las violaciones se fijen en la implementación, evitando que la deriva
        se convierta en un problema de mantenimiento costoso meses o años después
        \parencite{Klein2022}.

    \item \textbf{Modelo de Reflexión (Software Reflexion Models):} \textcite{Murphy2001}
        introdujeron el mecanismo formal para cuantificar la divergencia entre el
        grafo de dependencias del código fuente (\textit{source model}) y la
        arquitectura prescrita (\textit{high-level model}), clasificando cada
        relación en convergencias (permitidas), divergencias (violaciones =
        deriva) o ausencias (esperadas pero faltantes). Esta tripartición
        fundamenta directamente el motor de validación y la Tasa de Conformidad
        (CR) de este trabajo \parencite{Murphy2001}.
\end{itemize}